\section{Implementation Roadmap}
\label{sec:future:roadmap}

Phase Two completes the proposed system specified in
Chapter~\ref{chapter:solution}. The work is organised into nine phases,
labelled P1 through P9. Each phase has a single named output and a single
acceptance criterion that determines when the phase is complete. The total
effort budget is approximately twenty-one working days at one developer.
The roadmap introduces three abbreviations that recur in later sections:
the Automated Fare Collection (AFC) data source, the Global Positioning
System (GPS) data source, and the Quality of Service (QoS) reliability
framework adopted from Chapter~\ref{chapter:foundations}.

\begin{table}[!htbp]
  \centering
  \caption{Phase Two implementation roadmap}
  \label{tab:future:roadmap}
  \footnotesize
  \renewcommand{\arraystretch}{1.25}
  \begin{tabularx}{\linewidth}{p{0.7cm}p{4.6cm}p{1.5cm}X}
    \hline
    \textbf{Phase} & \textbf{Output} & \textbf{Effort} & \textbf{Acceptance criterion} \\
    \hline
    P1 & Warehouse Layers 2 to 4 (conformed dimensions, fact views, aggregate tables)
       & 2 days
       & The build command runs in under 30 seconds on the local sample and emits the
         ten dimension tables, three fact views, and eight aggregate tables; an
         integration check asserts row counts against the fact layer. \\
    P2 & First end-to-end theme (Demand and Crowding)
       & 3 days
       & All five charts in the theme render against the warehouse; the Multiple
         Coordinated Views model propagates selection across all five; the
         entity-detail panel opens on any route click. \\
    P3 & Spatial and Connectivity theme
       & 2 days
       & The hexagonal-bin map renders at three cell sizes; the per-stop demand view
         renders the 4{,}710 stops without user-interface lag. \\
    P4 & School-Aware Demand theme
       & 2 days
       & The schools-to-routes treemap renders; the school-routeness classifier
         renders for all routes; the school proximity map renders at district
         resolution. \\
    P5 & Financial and Operations plus Network Health themes
       & 3 days
       & The financial quadrant is interactive; the QoS-1, QoS-2, QoS-3 indicators
         are computed on the local GPS sample; the anomaly view supports the
         threshold control. \\
    P6 & Routes theme with time-distance diagram
       & 3 days
       & Cross-route and per-route views both render; the time-distance diagram
         supports zoom and trip hover. \\
    P7 & Cross-theme polish (overview index, entity-detail panel for all entity
         types, URL state)
       & 2 days
       & The overview index renders; the entity-detail panel works for all six
         entity types; the URL encodes the active theme and filters. \\
    P8 & Narrative scrollytelling layer
       & 2 days
       & The four planned narratives are scroll-driven; each scene deep-links to the
         corresponding theme view. \\
    P9 & Final polish and initial evaluation
       & 2 days
       & School geocoding is verified for at least 50 of the 64 schools; the full
         fifty-day GPS dataset is loaded if available; an end-to-end demonstration
         screencast is produced. \\
    \hline
  \end{tabularx}
\end{table}

\noindent\textbf{Content.} The plan content is the nine phases recorded in
Table~\ref{tab:future:roadmap}. Each phase produces one or two artefacts that
the next phase depends on. The acceptance criteria are verifiable by a single
command or a single demonstrable user interaction, which avoids open-ended
evaluation at phase boundaries and keeps the cost of declaring a phase
complete fixed.

\noindent\textbf{Time.} The total estimated effort is twenty-one working days
at one developer. The critical path is P1 followed by P2, which together
account for about five working days; the warehouse layers built in P1 are a
hard prerequisite for the first end-to-end theme exercised in P2. Phases
P3 and P4 may run in parallel after P2 completes because their underlying
aggregate tables are independent of one another. Phases P5 through P9 each
consume one or more aggregate tables produced in P1 and follow the order in
which the corresponding analytical themes depend on each other.

\noindent\textbf{Personnel.} Phase Two is single-author. Time is tracked per
phase rather than per task, which matches the granularity of the acceptance
criteria and keeps the tracking overhead small. Two advisor reviews are
scheduled: at the end of P2, after the first end-to-end theme has been
demonstrated, and at the end of P6, after all themes are rendering. The two
review points coincide with the two natural integration milestones in the
plan.

\noindent\textbf{Expected outcomes.} At the end of Phase Two, the proposed
system renders thirty out of thirty-one charts at sub-fifty-millisecond
warm latency, supports the Multiple Coordinated Views interaction within
every theme and the entity-detail interaction on every clickable element,
persists state to a shareable URL, and exports every chart in three formats
from a single one-click control. The remaining chart, the network-wide
arrival-time prediction, is treated as an extension in
Section~\ref{sec:future:extensions} because its underlying construction
depends on the full fifty-day GPS window rather than on the local sample.
